\section{Threat Model and Claim Boundary}
\label{sec:threat-model}

We study a honeypot whose attacker-facing output is generated by a provider
hosted model but accepted by a local, deterministic serving loop. The attacker
may observe and re-query a session, start fresh sessions, inject instructions
through commands or retrieved content, submit malformed or high-rate input,
and adapt to previously observed verifier behavior. Cross-session probing,
provider refusal, and retry pressure are in scope. A compromised host,
container escape, provider-side compromise, and operator collusion are outside
the in-system security property and are deployment prerequisites rather than
properties certified by the verifier.

The trusted computing base is the deterministic verifier, the fiction-graph
state and its isolation boundary, the canary registry, vetted fallback
transcripts, the gateway's key and retry policy, and the append-only audit
record. The LLM, retrieved CTI, attacker input, and provider availability are
untrusted or unreliable inputs. The serving claim is deliberately narrow: a
draft is served only when it passes the declared checks, otherwise the system
falls back or stops at the human gate.

\begin{table*}[t]
\centering
\caption{Security properties, evidence, and nonclaims. All evidence is bounded
to the named workload or declared command surface.}
\label{tab:threat-boundary}
\scriptsize
\setlength{\tabcolsep}{3pt}
\begin{tabular}{p{0.14\textwidth}p{0.20\textwidth}p{0.29\textwidth}p{0.29\textwidth}}
\toprule
Property & Attacker pressure & Evidence & Explicit nonclaim \\
\midrule
Fail-closed serving & Re-query, malformed draft, identity or consistency
drift & 105-case verifier suite; 778 induced drafts from three author families
of which the serving path rejected 554 of the 572 violating drafts and
admitted 18, refusing 17 of 206 benign; 45 graph re-queries; boundary probes
stopping 4 of 7 authored disclosures and 5 of 10 credential formats &
Not open-world resistance, realism, or coverage against arbitrary adaptive
drafts. A fluent, on-fiction secret disclosure is not detected in general: the
gate covers the credential formats it was written against, and the probes'
residual risk is named in \S\ref{sec:evaluation}. \\
Supported claims & Unsupported answers on the eight-command declared surface
(\texttt{whoami}, \texttt{id}, \texttt{hostname}, \texttt{uname},
\texttt{echo}, \texttt{uptime}, \texttt{date}, \texttt{cat} of
\texttt{/etc/os-release}) & 123 supported-claim cases over exactly that
surface & Closed-surface recall only; no general semantic guarantee, and no
demonstrated prompt-injection or retrieval-poisoning resistance. \\
State isolation & Cross-session or cross-host drift & Host-fiction and graph
consistency artifacts, including 4/4 deterministic drift cases & No
multi-tenant or compromised-host guarantee. \\
Availability & Provider refusal, timeout, quota, and retry pressure &
Deterministic fallback and offline failure-path battery & No field-load
availability, cost, or tail-latency claim. \\
Grounding & Retrieval removed from the same fixed scenario set & Archived
model-coherent analysis-channel comparison and paired bootstrap artifact &
Not a clean served-text causal result; fresh provider reproduction remains
external work. \\
\bottomrule
\end{tabular}
\end{table*}

The held-out rescue benchmark and the adaptive ladder are treated as
exploratory controller diagnostics. The rescue risk floor is derived from the
controller itself, and the ladder's five-pair replay has no power to establish
attacker engagement or security benefit.

Two attacker capabilities named above are declared in scope but
unevaluated, and we state the consequence rather than the aspiration.
Instruction injection through commands or retrieved content is represented in
the induced-draft battery by two cases, both of which carry targeted-scoring
misses, with one variant of one case passing the gate; the verifier is
therefore not credited with injection resistance anywhere in this paper.
Retrieval poisoning is not evaluated at all. The claim boundary in
Table~\ref{tab:threat-boundary} is the wording the evidence supports; stronger
claims require independently constructed attacks, independent quality
evaluation, or field deployment.
